\section{Protocolo de validação do Text-to-SQL}

A avaliação do subsistema Text-to-SQL foi conduzida por \textit{pipeline} automatizado desacoplado do canal WhatsApp, garantindo reprodutibilidade (detalhes em Apêndice~\ref{ap:artefatos}). A validação do subsistema RAG segue protocolo complementar descrito na Seção seguinte.

\section{Protocolo de validação do RAG}

Complementarmente ao Text-to-SQL, conduziu-se avaliação da recuperação documental sobre o Diário Oficial municipal indexado em Qdrant (9\,346 \textit{chunks}, embedding \texttt{text-embedding-3-small}).

\subsection{Conjunto de referência RAG}

O \textbf{conjunto de referência RAG ampliado} contém $N_{\text{RAG}} = 40$ perguntas estratificadas em dois estratos:

\begin{itemize}
    \item \textbf{com\_numero} ($n = 32$): decretos sorteados uniformemente do inventário do índice (seed 42), com pergunta-template contendo número e ano explícitos do ato;
    \item \textbf{tematico} ($n = 8$): perguntas curadas sobre conteúdo normativo, reaproveitadas da amostra temática inicial, sem depender exclusivamente de número na formulação.
\end{itemize}

A inclusão dos itens \textit{com\_numero} é independente do \textit{rank} de recuperação --- critério de amostra honesta. Cada item possui documento de referência validado no índice (cabeçalho do ato presente em \textit{page\_content}).

\subsection{Indicadores de recuperação}

Para cada pergunta $i$, registra-se o ranking do documento esperado entre os $k$ primeiros \textit{chunks} retornados ($k = 5$):

\begin{equation}
\text{Recall@}k = \frac{1}{N_{\text{RAG}}} \sum_{i=1}^{N_{\text{RAG}}} \mathbb{1}\!\left[\text{doc esperado} \in \text{top-}k_i\right]
\end{equation}

A \textbf{taxa de recuperação média} (MRR --- \textit{Mean Reciprocal Rank}) mede a posição média do primeiro acerto:

\begin{equation}
\text{MRR} = \frac{1}{N_{\text{RAG}}} \sum_{i=1}^{N_{\text{RAG}}} \frac{1}{\text{rank}_i}
\end{equation}

onde $\text{rank}_i$ é a posição do \textit{chunk} que contém o ato esperado, ou zero se ausente do top-$k$. Reportam-se também intervalos de confiança de Wilson (95\%) \cite{Agresti1998} para as proporções de acerto.

\subsection{Escopo e limitações do protocolo RAG}

A avaliação compara estratégias de recuperação (busca densa, deduplicação, filtro por tipo de ato e busca híbrida por número/ano) sobre o mesmo \textit{golden}. Não avalia geração da resposta final em linguagem natural nem o roteamento do orquestrador.


\subsection{Notação e pré-condições}

Para cada consulta $i$ no conjunto de referência ($N = 80$):

\begin{itemize}
    \item $q_i$ --- pergunta em linguagem natural;
    \item $\text{SQL}_i^{(r)}$ --- consulta de referência (gabarito);
    \item $\text{SQL}_i^{(g)}$ --- consulta gerada pelo modelo;
    \item $R_i$ --- resultado tabular de $\text{SQL}_i^{(r)}$;
    \item $G_i$ --- resultado tabular de $\text{SQL}_i^{(g)}$;
    \item $C_i$ --- conjunto de colunas de resposta declaradas no gabarito.
\end{itemize}

Define-se a função de normalização de célula $\nu(\cdot)$, aplicada antes de qualquer comparação: nulos unificados; números arredondados em quatro casas decimais; strings sem espaços laterais e em maiúsculas; strings numéricas convertidas com vírgula substituída por ponto.

\subsection{Taxa de SQL válida (VSR)}

A \textbf{taxa de SQL válida} (\textit{Valid SQL Rate}, VSR) mede se a consulta gerada é sintaticamente correta e executável sobre o esquema municipal:

\begin{equation}
\text{VSR} = \frac{1}{N} \sum_{i=1}^{N} \mathbb{1}\!\left[\text{SQL}_i^{(g)} \text{ executa sem erro}\right]
\end{equation}

onde $\mathbb{1}[\cdot]$ é a função indicadora. Quando a SQL gerada falha, os indicadores de equivalência de resultado para a consulta $i$ são nulos.

\subsection{Equivalência funcional (EF)}

A \textbf{equivalência funcional} (EF) verifica se o cidadão receberia as grandezas corretas, independentemente de aliases de coluna ou colunas auxiliares no resultado gerado. Trata-se da métrica adotada neste trabalho para comprovar a utilidade do resultado tabular.

Diferentemente do \textit{Execution Accuracy} (EX) usado em Spider e BIRD --- que exige correspondência estrita de nomes de colunas e do multiconjunto de pares (coluna, valor) ---, a EF restringe a comparação às colunas de resposta declaradas no conjunto de referência, admite mapeamento entre aliases de coluna e aplica tolerância numérica $\tau$ em grandezas fiscais. Essa definição penaliza menos formulações SQL semanticamente equivalentes, alinhando-se ao objetivo de comprovar utilidade para o cidadão.

Para cada consulta $i$ com VSR positivo, exige-se $|R_i| = |G_i|$ (mesma cardinalidade de linhas). Para cada coluna $k \in C_i$, deve existir alguma coluna $j$ em $G_i$ cujo vetor de valores normalizados ao longo das linhas seja equivalente ao vetor da coluna correspondente em $R_i$, com tolerância numérica $\tau = 0{,}5\%$ em valores fiscais:

\begin{equation}
\left|\nu(a) - \nu(b)\right| / \max\!\left(|\nu(b)|, \varepsilon\right) \leq \tau, \quad \tau = 0{,}005
\end{equation}

onde $\varepsilon$ evita divisão por zero. A equivalência funcional agregada é:

\begin{equation}
\text{EF} = \frac{1}{N} \sum_{i=1}^{N} \mathbb{1}\!\left[\text{EF}_i = 1\right]
\end{equation}

com $\text{EF}_i = 1$ quando todas as colunas em $C_i$ possuem correspondente equivalente em $G_i$ sob as regras acima. Proporções agregadas são reportadas com intervalos de confiança de Wilson (95\%) \cite{Agresti1998}; comparações pareadas E1 \textit{versus} Prompt B utilizam teste de McNemar quando aplicável.

\subsection{Indicadores complementares}

\begin{description}
    \item[NEA, TSA, CHS] Indicadores auxiliares de resultado não vazio, seleção de tabelas e presença de filtros relevantes.
    \item[JAR --- \textit{Judge Agreement Rate}] Proporção de respostas em linguagem natural classificadas como corretas ou parcialmente corretas por LLM-as-Judge \cite{Zheng2023LLMJudge}, em amostra estratificada ($M = 30$):
    \begin{equation}
    \text{JAR} = \frac{1}{M} \sum_{j=1}^{M} \mathbb{1}\!\left[\text{veredito}_j \in \{\text{correto}, \text{parcialmente\_correto}\}\right]
    \end{equation}
\end{description}

A implementação completa, incluindo variantes diagnósticas, está documentada no repositório público de artefatos do trabalho (Apêndice~\ref{ap:artefatos}).

\subsection{Organização dos lotes experimentais}

A validação empírica organizou-se em três lotes sobre dados fixos (conjunto SQL de 80 casos; conjunto RAG ampliado de 40 perguntas):

\begin{enumerate}
    \item \textbf{Lote I} --- comparação de modelo e tamanho de contexto (E1--E4), fixando GPT-4.1-mini e Prompt A como referência;
    \item \textbf{Lote II} --- ablation de engenharia de \textit{prompt} (Prompt B), mantendo modelo e conjunto de referência do Lote I;
    \item \textbf{Lote III} --- comparação de estratégias de recuperação documental (busca densa, deduplicação, filtro por tipo de ato e busca híbrida), mantendo índice e embedding fixos.
\end{enumerate}

Cada lote isolou uma variável experimental.

\subsection{Experimentos do Lote I}

\begin{enumerate}
    \item \textbf{E1 --- Baseline}: GPT-4.1-mini com Prompt A (compacto);
    \item \textbf{E2 --- Ablação de contexto}: GPT-4.1-mini com \textit{prompt} estendido;
    \item \textbf{E3 --- Ablação de modelo (econômico)}: GPT-4o-mini com Prompt A;
    \item \textbf{E4 --- Ablação de modelo (avançado)}: GPT-4o com Prompt A.
\end{enumerate}

\section{Ambiente experimental}

\begin{itemize}
    \item DuckDB local sobre bases tabulares SICOM/CGU consolidadas;
    \item Modelos OpenAI via API oficial, temperatura 0,0;
    \item Saída estruturada em JSON contendo a consulta SQL;
    \item Persistência incremental em arquivo de \textit{checkpoint} para retomada de experimentos.
\end{itemize}
